\documentclass[10pt,a4paper]{article}
\usepackage{fontspec}
\setmainfont{Liberation Sans}
\usepackage{polyglossia}
\setdefaultlanguage{polish}
\usepackage[a4paper,margin=19mm,headheight=15pt]{geometry}
\usepackage{xcolor,tabularx,booktabs,array,enumitem,fancyhdr,titlesec,hyperref}
\definecolor{navy}{HTML}{263650}
\definecolor{blue}{HTML}{5475A5}
\definecolor{pale}{HTML}{EFF3F8}
\definecolor{muted}{HTML}{5D6878}
\hypersetup{colorlinks=true,urlcolor=blue}
\pagestyle{fancy}
\fancyhf{}
\fancyhead[L]{\small\color{muted}Tutorland}
\fancyhead[R]{\small\color{muted}Stan prac}
\fancyfoot[C]{\small\color{muted}\thepage}
\renewcommand{\headrulewidth}{0.3pt}
\setlength{\parindent}{0pt}
\setlength{\parskip}{5pt}
\setlength{\emergencystretch}{3em}
\sloppy
\setlist[itemize]{leftmargin=*,itemsep=3pt,topsep=3pt}
\titleformat{\section}{\Large\bfseries\color{navy}}{\thesection.}{0.55em}{}
\newcolumntype{Y}{>{\raggedright\arraybackslash}X}
\begin{document}
\begin{titlepage}
\thispagestyle{empty}
\vspace*{32mm}
{\large\color{blue}\textbf{RAPORT POSTĘPU · 26 PAŹDZIERNIKA 2026}}
\vspace{12mm}
{\Huge\bfseries\color{navy}Tutorland\par}
\vspace{4mm}
{\LARGE Aktualny stan aplikacji\par}
\vspace{13mm}
\color{blue}\rule{\textwidth}{2pt}
\vspace{9mm}
{\large Przegląd aplikacji, lokalnego API, przygotowania PHP/MySQL i planowanego hostingu.}
\vfill
\colorbox{pale}{\parbox{\dimexpr\textwidth-2\fboxsep\relax}{\textbf{W skrócie}\\
Działa katalog publiczny, logowanie i rejestracja z wyborem roli, profile, rezerwacje, panele ucznia i korepetytora oraz opinie. Lokalny tryb developerski korzysta z Node.js i CSV; dla hostingu PHP przygotowano API PHP, schemat MySQL i konfigurację wdrożenia. Zoom OAuth wymaga ustawień Zoom, sekretów na hostingu oraz potwierdzenia, że hosting dopuszcza połączenia wychodzące.}}\par
\vspace{12mm}
{\small\color{muted}Raport opisuje kod projektu. Nie oznacza wdrożenia publicznego serwera ani aktywnego połączenia z zewnętrznym kontem Zoom.}
\end{titlepage}

\section{Zaimplementowane funkcje}
\begin{itemize}
\item \textbf{Katalog publiczny:} karty i profile korepetytorów, wyszukiwanie po imieniu i nazwisku, filtr przedmiotu i poziomu nauczania oraz podgląd dostępnych terminów przez gości.
\item \textbf{Konta i role:} osobne rejestracja i logowanie z wyborem ucznia lub korepetytora, sesja w ciasteczku \texttt{HttpOnly}, hasła przechowywane jako hashe scrypt oraz wylogowanie. Rola jest sprawdzana po stronie API.
\item \textbf{Tryb developerski:} lokalne API Node obsługuje konta, katalog, profile, dostępność, rezerwacje i opinie; zapisuje dane w prywatnych plikach CSV w \texttt{server/data/}.
\item \textbf{Tryb wdrożeniowy:} dodano API PHP w \texttt{public/api/index.php} z tymi samymi ścieżkami \texttt{/api/...}, obsługą kont, sesji, ról, profili, dostępności, rezerwacji, opinii i Zoom OAuth.
\item \textbf{MySQL:} utworzono \texttt{tutorland-database.sql} poza katalogiem aplikacji. Schemat obejmuje konta, profile, przedmioty, poziomy, dostępność, dni wolne, rezerwacje, opinie, sesje, stany OAuth i szyfrowane tokeny Zoom.
\item \textbf{Rezerwacje:} uczeń może zapisać się na dostępny termin lekcji 40-minutowej. API weryfikuje rolę, przedmiot, dostępność, dni wolne, przeszłe godziny i nakładające się rezerwacje.
\item \textbf{Panel ucznia:} lista zaplanowanych i zakończonych spotkań, linki do spotkań Zoom, historia własnych opinii oraz formularze dodawania, edycji i usuwania. Opinię można dodać po zakończonej lekcji; obowiązuje jedna aktywna opinia na ucznia i korepetytora.
\item \textbf{Panel korepetytora:} edycja profilu, przedmiotów, poziomów i ceny; suwaki do ustawiania godzin dostępności w tygodniu; kalendarz liczby rezerwacji w bieżącym miesiącu i oznaczanie dni wolnych w kolejnym miesiącu.
\item \textbf{Powiadomienia korepetytora:} zaplanowane i zakończone lekcje oraz opinie uczniów. Korepetytor może je przeglądać, ale nie edytować.
\item \textbf{Połączenie Zoom:} OAuth korepetytora, pobieranie zdjęcia profilu Zoom i spotkania 40-minutowe. Tokeny są szyfrowane lokalnie w CSV-backendzie lub w tabeli MySQL w backendzie PHP. Uczeń otrzymuje link dołączenia, a link prowadzący jest udostępniany korepetytorowi.
\item \textbf{Dane konta demonstracyjnego:} usunięto wskazane konto testowe z aplikacji i źródłowych danych demonstracyjnych. Uczniowie nie mają zdjęć profilowych. Nazwę poziomu zmieniono na „Szkoła Podstawowa”.
\end{itemize}

\section{Konfiguracja i ograniczenia}
\small
\begin{tabularx}{\textwidth}{>{\bfseries}p{35mm}Y}
\toprule
Obszar & Stan i wymagane działanie\\
\midrule
Lokalny start & \texttt{npm run dev} uruchamia Vite i API Node. Zmienne dane developerskie trafiają do ignorowanego katalogu \texttt{server/data/}.\\
Produkcja & Przygotowano statyczny build React z API PHP oraz routingiem Apache przez pliki \texttt{.htaccess}. Na hosting należy wgrać zawartość \texttt{dist/}, a nie \texttt{node\_modules} ani całe repozytorium.\\
MySQL & Zaimportuj \texttt{tutorland-database.sql} do bazy utworzonej w panelu/phpMyAdmin. Konfiguracja połączenia jest tworzona osobno jako ignorowany plik \texttt{api/config.local.php}; wartości z panelu hostingu nie trafiają do repozytorium.\\
Zoom OAuth & Przygotowany adres callbacku to \texttt{https://tutorland.infinityfreeapp.com/api/zoom/callback}; lista dozwolonych obejmuje domenę \texttt{tutorland.infinityfreeapp.com}. Potrzebne są zakresy odczytu profilu i tworzenia spotkań oraz prywatna konfiguracja Zoom na hostingu.\\
Linki spotkań & Rezerwacja bez połączonego Zoom pozostaje oczekująca; po połączeniu serwer próbuje utworzyć spotkania oczekujące. Błędy Zoom są zapisywane i pokazywane w statusie rezerwacji.\\
Poczta e-mail & Aplikacja nie wysyła e-maili; powiadomienia są dostępne w panelach kont.\\
Środowisko & Kod jest przygotowany do wdrożenia, ale nie został jeszcze wgrany. HTTPS, dane połączenia MySQL, klucz szyfrowania, backupy i rzeczywisty przepływ Zoom wymagają konfiguracji na hostingu.\\
Ograniczenie hostingu & Serwer musi mieć PHP cURL i dopuszczać połączenia HTTPS wychodzące do Zoom. To trzeba potwierdzić testem po publikacji.\\
\bottomrule
\end{tabularx}
\normalsize

\section{Sprawdzenia}
\begin{itemize}
\item Testy Node sprawdzają kodowanie CSV oraz podstawowe przepływy API: konta, sesje, dostęp do profilu, zapis rezerwacji, odrzucanie duplikatów i opinii przed odbytą lekcją.
\item \texttt{npm run lint} przechodzi bez ostrzeżeń.
\item \texttt{npm run build} generuje frontend produkcyjny.
\item PHP API przechodzi kontrolę składni PHP 8.3. Schemat SQL został zaimportowany do tymczasowej bazy MariaDB; test integracyjny sprawdził rejestrację, logowanie, profile, dostępność, kolizję rezerwacji i cykl opinii.
\item Nie wykonano jeszcze wdrożenia na InfinityFree ani rzeczywistego połączenia OAuth/Zoom.
\end{itemize}

\end{document}
